\Article{297}
Wer den vorigen Beweis genauer betrachtet, erkennt leicht, dass
die Fälle I und II ohne Einwand erledigt sind. Die Beweise der
übrigen Fälle stützen sich aber auf die Existenz von Hilfszahlen.
Solange diese nicht bewiesen ist, verliert die Methode offenbar
jede Beweiskraft. Diese Annahmen mögen so einleuchtend sein, dass
sie einem weniger aufmerksamen Leser gar keines Beweises zu bedürfen
scheinen, und sie heben den zu beweisenden Satz gewiss auf den
höchsten Grad der \emph{Wahrscheinlichkeit}; verlangt man aber
geometrische Strenge, so darf man sie keinesfalls ohne Beweis
zulassen. Was die Annahme in IV und V betrifft, es gebe eine Primzahl
\(r\) der Form \(4n+1\), die Nichtrest zweier anderer gegebener
Primzahlen \(p,q\) ist, so folgt aus Abschnitt IV leicht, dass
alle Zahlen unter \(4pq\), die dazu teilerfremd sind (ihre Anzahl
ist \(2(p-1)(q-1)\)), gleichmäßig auf vier Klassen verteilt sind. Eine
enthält Nichtreste von beiden \(p,q\), die übrigen drei Reste
von \(p\) und Nichtreste von \(q\), Nichtreste von \(p\)
und Reste von \(q\) sowie Reste von beiden \(p,q\). In jeder
Klasse ist die Hälfte der Zahlen von der Form \(4n+1\), die andere
von der Form \(4n+3\). Es gibt darunter also \(\frac14(p-1)(q-1)\) Nichtreste
von beiden \(p,q\) der Form \(4n+1\); diese seien \(g.g',g''\) usw.
Die übrigen \(\frac74(p-1)(q-1)\) Zahlen seien \(h,h',h''\) usw.\ Offenbar sind
alle in den Formen \(4pqt+g,4pqt+g',4pqt+g''\) usw.\ \((G)\) enthaltenen Zahlen
ebenfalls Nichtreste von \(p,q\) der Form \(4n+1\). Um die Annahme
zu begründen, braucht man daher nur zu beweisen, dass unter den
Formen \((G)\) sicher \emph{Primzahlen} vorkommen. Das erscheint
schon an sich sehr plausibel, da diese Formen zusammen mit
\(4pqt+h,4pqt+h'\) usw.
\SourceLeaf{362}
\((H)\) alle zu \(4pq\) teilerfremden Zahlen und somit auch
alle Primzahlen (außer \(2,p,q\)) umfassen. Es scheint kein Grund
vorzuliegen, weshalb die Reihe der Primzahlen nicht gleichmäßig
auf jene Formen verteilt sein sollte, sodass ein Achtel auf
\((G)\), die übrigen auf \((H)\) entfallen. Eine solche Überlegung
ist aber offenbar weit von geometrischer Strenge entfernt.
Der berühmte Le Gendre selbst räumt ein, dass der Beweis des
Satzes, unter einer solchen Form \(kt+l\) mit gegebenen teilerfremden
Zahlen \(k,l\) und veränderlichem \(t\) kämen sicher Primzahlen
vor, ziemlich schwierig erscheine. Er deutet beiläufig eine
Methode an, die vielleicht dorthin führen könnte. Uns scheinen
aber viele Voruntersuchungen nötig, bevor man wenigstens auf
diesem Weg zu einem strengen Beweis gelangen kann. — — Zur
anderen Annahme (III, zweite Methode), es gebe eine Primzahl
\(r\) der Form \(3n+3\), von der eine andere gegebene Primzahl
\(p\) der Form \(4n+1\) Nichtrest ist, hat der berühmte Le Gendre
überhaupt nichts hinzugefügt. Wir haben oben in Art.\ 129
bewiesen, dass es sicher Primzahlen gibt, von denen \(p\)
Nichtrest ist. Unsere Methode scheint aber kaum geeignet,
die Existenz solcher Primzahlen nachzuweisen, \textit{die
zugleich von der Form} \(4n+3\) sind, wie es hier, aber nicht
in unserem ersten Beweis nötig ist. Die Wahrheit dieser
Annahme können wir allerdings so leicht beweisen: Nach
Art.\ 287 gibt es ein positives Geschlecht binärer Formen
mit der Determinante \(-p\) und dem Charakter \(3,4;\ Np\). Sei
\((a,b,c)\) eine solche Form und \(a\) ungerade, was wir voraussetzen
dürfen. Dann ist \(a\) von der Form \(4n+3\) und entweder
selbst eine Primzahl oder hat wenigstens einen Primfaktor
\(r\) der Form \(4n+3\). Es ist aber \(-pRa\), also auch
\(-pRr\), woraus \(pNr\) folgt. Dabei muss man sorgfältig
beachten, dass sich die Sätze aus Art.\ 263, 287 auf den
Fundamentalsatz stützen. Es wäre also ein Zirkelschluss,
einen Teil dieses Satzes darauf zu gründen. — — Schließlich
ist die Annahme in III der ersten Methode noch viel weniger
begründet, sodass wir hier nichts weiter darüber hinzuzufügen
brauchen.

Wir fügen noch eine Bemerkung zu Fall V hinzu, der durch
die vorige Methode zwar nicht ausreichend bewiesen, durch
die folgende aber bequem erledigt wird. Wäre dort zugleich
\(pNq,qNp\), so wäre \(-pRq,-qRp\). Daraus folgt leicht, dass \(-1\)
eine charakteristische Zahl der Form \((p,0,q)\) ist, die daher
nach der Theorie der ternären Formen durch \(xx+yy+zz\) dargestellt
werden kann. Es sei
\[ptt+quu=(\alpha t+\beta u)^2+(\alpha't+\beta'u)^2+(\alpha''t+\beta''u)^2\]
oder
\[\begin{gathered}\alpha\alpha+\alpha'\alpha'+\alpha''\alpha''=p,\\\beta\beta+\beta'\beta'+\beta''\beta''=q,\\\alpha\beta+\alpha'\beta'+\alpha''\beta''=0\end{gathered}\]

\SourceLeaf{363}
Dann sind nach den Gleichungen 1 und 2 alle \(\alpha,\alpha',\alpha'',\beta,\beta',\beta''\) ungerade;
die dritte Gleichung kann aber dann offenbar nicht bestehen.
— Auf ganz ähnliche Weise lässt sich auch Fall II erledigen.
